\documentclass[10pt]{article}
\usepackage[a4paper,margin=1in]{geometry}
\usepackage{amsmath,amssymb,amsthm}
\usepackage{booktabs}
\usepackage{enumitem}
\usepackage[hidelinks]{hyperref}
\newcommand{\tool}{\textsc{MiniTuner}}
\newtheorem{theorem}{Theorem}

\title{\textsc{MiniTuner}: Exact Candidate-Guided Pareto Search\\
for Floating-Point Function Implementations}
\author{Your Name\\Your Institution}
\date{}

\begin{document}
\maketitle

\begin{abstract}
Numerical libraries offer elementary-function implementations with different
costs, valid domains, and error bounds. Choosing one implementation at every
call site is a discrete multi-objective problem. \tool constructs an exact
rational cost--error model, obtains a small set of feasible candidate solutions
directly from that model, and uses their certified objective values to add sound
dominance cuts. Z3 enumerates the complete remaining Pareto frontier, while
FPTaylor supplies symbolic error coefficients and validates every reported
configuration. No statistical model participates in the
production pipeline. The FPCore front end preserves source domains and attempts
conservative binary64-valid domain inference when bounds are missing. Every
benchmark produces a frontier image or an explicit status image, together with
phase-separated timing data.
\end{abstract}

\section{Problem Formulation}

Let $S=\{s_1,\ldots,s_n\}$ be the tunable function sites in an expression.
Site $s_i$ has implementations $I_i=\{I_{i1},\ldots,I_{im_i}\}$. An
implementation has measured cost $c_{ij}$, valid domain $D_{ij}$, and local
error specification
\begin{equation}
 |\widehat f_{ij}(x)-f_i(x)| \leq
 \epsilon_{ij}|f_i(x)|+\delta_{ij},\qquad x\in D_{ij}.
\end{equation}
A configuration $q=(j_1,\ldots,j_n)$ selects one implementation per site.
FPTaylor derives conservative whole-expression coefficients, giving the exact
selector objectives
\begin{align}
 C(q)&=C_{\mathrm{fixed}}+\sum_i c_{ij_i},\\
 E(q)&=E_{\mathrm{base}}+
 \sum_i(A_i\epsilon_{ij_i}+B_i\delta_{ij_i}).
\end{align}
The Z3 encoding retains rational coefficients. A pair $(C_1,E_1)$ dominates
$(C_2,E_2)$ when both objectives are no worse and at least one is strictly
better. The output contains every distinct nondominated rational objective
pair.

\section{Exact Candidate Guidance}

After domain filtering, tuple deduplication, and safe local dominance
reduction, each site is encoded by exactly-one Boolean choices. \tool obtains
candidate solutions from this same exact model:
\begin{enumerate}[leftmargin=*]
 \item lexicographically minimum cost and then error;
 \item lexicographically minimum error and then cost; and
 \item additional normalized positive cost--error scalarizations.
\end{enumerate}
Each optimizer result is already a feasible exact configuration. For a
certified point $(C_s,E_s)$, \tool adds
\begin{equation}
 C<C_s\;\lor\;E<E_s\;\lor\;(C=C_s\land E=E_s).
 \label{eq:cut}
\end{equation}

\begin{theorem}[Dominance-cut safety]
Equation~\eqref{eq:cut} removes no distinct nondominated objective pair.
\end{theorem}
\begin{proof}
A removed feasible pair has $C\geq C_s$ and $E\geq E_s$ and is not equal to
the stored candidate pair. The certified candidate therefore weakly improves
both objectives and strictly improves at least one, so the removed pair is
dominated.
\end{proof}

The candidate count is configurable. Setting it to zero gives the unguided
baseline. Guided and unguided runs are required to produce identical rational
frontiers; candidate discovery is an acceleration heuristic, not a change to
the feasible set.

\section{Exact Enumeration and Validation}

The default engine uses Z3's Pareto optimization. A staircase engine is kept
for comparison. Optional disjoint cost bands use independent solver processes;
all bands must complete before their union is labelled complete. Band outputs
are merged by exact rational objective pairs and reduced once more for global
dominance.

Every enumerated configuration is then fixed and checked by FPTaylor over the
complete analyzed domain. Failed validation is retried by bounded domain
subdivision. A timeout, unknown solver result, unresolved domain, or failed
validation cannot produce a complete-frontier label.

\section{FPCore Domains}

Source finite bounds are preserved. Extra relational constraints may be
conservatively over-approximated by their enclosing box and are labelled as
such. When a variable lacks finite bounds, inference begins with the finite
binary64 range and applies operation-validity restrictions for square root,
logarithms, exponentials, division, powers, tangent, and arithmetic overflow.
Only regions proved finite and operation-valid are admitted. Disconnected or
unresolved domains are reported explicitly rather than replaced with an
arbitrary experimental interval.

Fixed-only expressions have no choice Boolean and produce a one-point
frontier. Unsupported precision or control flow produces a status result. Thus
every suite entry is accounted for without presenting a changed program as the
original benchmark.

\section{Artifacts and Reproduction}

The complete suite command is
\begin{verbatim}
nix develop path:. --command ./run_all_fpcore.sh \
  --timeout 120 --jobs 2 --guidance-points 4
\end{verbatim}
All generated material is under \texttt{output/fpcore}: imported inputs,
frontier CSVs, images, logs, resumable checkpoints, \texttt{summary.csv},
\texttt{timings.csv}, and \texttt{run.json}. Each benchmark has an image. A
successful image shows cost versus certified error; an unsuccessful image
shows the exact status and reason.

Timing is separated into import/domain handling, exact preprocessing,
candidate certification, Z3 enumeration, FPTaylor validation, plotting, and
end-to-end wall time. This separation is necessary because a cut may reduce
enumeration time while candidate discovery still increases total time on a
small instance.

\section{Evaluation Methodology}

Correctness tests compare guided and zero-candidate frontiers using exact
rational pairs. Small instances are also exhaustively enumerated. Domain tests
cover binary64 boundaries, invalid square roots and logarithms, division by
zero, negative powers, exponential overflow, and disconnected domains. Suite
tests require one summary row, timing row, and image per FPCore entry, and
verify that partial runs never overwrite complete results.

Performance evaluation reports candidate-certification and Z3 times
separately. A candidate mechanism is worthwhile only when its end-to-end saving
exceeds its exact pre-solve cost; no universal speedup is assumed.

\section{Limitations}

Exact real-error analysis cannot assign a finite error to executions producing
NaN or infinity. Domain inference is therefore conservative and may reject a
benchmark when it cannot prove complete valid coverage within its resource
budget. General control flow, unresolved relational domains, and unsupported
precisions remain explicit status outcomes. Exact completeness may still be
expensive because sound candidate cuts cannot remove any potentially
nondominated region.

\section{Conclusion}

\tool uses exact candidate solutions to accelerate an exact Pareto search.
Candidate configurations, dominance cuts, enumeration, objective merging, and
final validation all preserve rational frontier semantics. The central
invariant is that only a feasible exact point may exclude its strictly
dominated region. Organized per-benchmark outputs and explicit failure images
make both correctness and runtime visible across the FPCore suite.

\begin{thebibliography}{9}
\bibitem{fptaylor}
A.~Solovyev, C.~Jacobsen, Z.~Rakamari\'c, and G.~Gopalakrishnan.
Rigorous Estimation of Floating-Point Round-off Errors with Symbolic Taylor
Expansions. In \emph{FM}, 2015.
\bibitem{fpbench}
N.~Damouche et al. Toward a Standard Benchmark Format and Suite for
Floating-Point Analysis. In \emph{NSV}, 2016.
\bibitem{z3}
L.~de Moura and N.~Bj\o{}rner. Z3: An Efficient SMT Solver. In \emph{TACAS},
2008.
\end{thebibliography}
\end{document}
